\chapter*{摘要}
\addcontentsline{toc}{chapter}{摘要}
课程设计包括两项图像分类、一项眼底关键点定位与一项像素分割。植物幼苗分类使用
分层五折训练的 ConvNeXt 和 EfficientNet。对前景裁剪中原图坐标与缩放坐标错用的
问题进行诊断并修复后，四配置集成的折外准确率记录为 0.9600，Kaggle 提交
\code{submission_v16_4model.csv} 的公开榜与私有榜均为 0.9735。面部表情分类在
六类 $48\times48$ 灰度图上比较输入尺度、骨干及融合策略。包含 ConvNeXt-Small
224 像素模型的集成与概率空间 bagging 的榜单准确率均为 0.7470。

眼底中心凹定位以 EfficientNet-B1 和 U 形解码器预测中心凹热图，将血管图作为
附加输入，并按本地模型与 VascX 预测的距离融合坐标。最终提交的公开榜与私有榜
坐标 RMSE 分别为 9.90927 和 19.54239 像素。DRIVE 视网膜血管分割使用 FR-UNet
公开权重，通过双阈值迭代、视野掩膜和连通域过滤生成提交掩膜，公开榜和私有榜
Dice 分别为 0.81766 和 0.82534。独立训练的原生分辨率 U-Net 达到 0.80780 的本地
整图验证 Dice，其对应公开榜表现低于 FR-UNet。

报告区分配置中的计划训练、日志中的实际训练及提交所用模型，并以代码、折外记录和
Kaggle 提交记录追溯图表。实验分析着重讨论预处理缺陷、验证与榜单分数的变化、
少量困难样本对 RMSE 的影响，以及视野与阈值口径对分割对照的影响。

\medskip
\noindent\textbf{关键词}\quad 迁移学习、折外验证、中心凹定位、视网膜血管分割、实验复现
